% ==============================================================================
% TRANG BÌA CHÍNH (TITLE PAGE)
% ==============================================================================
\begin{titlepage}
\begin{center}
    \textbf{\large BỘ GIÁO DỤC VÀ ĐÀO TẠO}\\[0.2cm]
    \textbf{\large TRƯỜNG ĐẠI HỌC GIAO THÔNG VẬN TẢI}\\[0.2cm]
    \textbf{\large KHOA CÔNG NGHỆ THÔNG TIN}\\[0.2cm]
    \textbf{\large BỘ MÔN KHOA HỌC MÁY TÍNH}\\[1.5cm]

    \IfFileExists{figures/utc_logo.png}{\includegraphics[width=3.5cm]{figures/utc_logo.png}\\[1.5cm]}{\vspace{2.5cm}}

    \textbf{\Large BÁO CÁO BÀI TẬP LỚN}\\[0.3cm]
    \textbf{\large HỌC PHẦN: XỬ LÝ ẢNH (IMAGE PROCESSING)}\\[1.0cm]

    \textbf{\LARGE ĐỀ TÀI SỐ 18:}\\[0.4cm]
    \textbf{\huge PHÁT HIỆN VẾT DẦU TRÀN TRÊN BIỂN TỪ ẢNH SAR SENTINEL-1}\\[0.3cm]
    \textbf{\Large Ứng dụng giám sát và cảnh báo sớm bờ biển miền Trung Việt Nam}\\[1.5cm]

    \begin{tabular}{ll}
        \textbf{Giảng viên hướng dẫn:} & TS. Phan Thanh Toàn \\
        \textbf{Nhóm sinh viên thực hiện:} & Nhóm X (Lớp CNTT - K63) \\
        \textbf{Trưởng nhóm (PM - TV1):} & [Họ và tên SV1 - Mã SV] \\
        \textbf{Thành viên lý thuyết (TV2):} & [Họ và tên SV2 - Mã SV] \\
        \textbf{Thành viên thuật toán 1 (TV3):} & [Họ và tên SV3 - Mã SV] \\
        \textbf{Thành viên thuật toán 2 (TV4):} & [Họ và tên SV4 - Mã SV] \\
        \textbf{Thành viên thực nghiệm (TV5):} & [Họ và tên SV5 - Mã SV] \\
    \end{tabular}

    \vfill
    \textbf{HÀ NỘI, NĂM 2026}
\end{center}
\end{titlepage}

% ==============================================================================
% TRANG LỜI CAM ĐOAN & LỜI CẢM ƠN
% ==============================================================================
\newpage
\pagenumbering{roman}
\setcounter{page}{1}

\section*{LỜI CAM ĐOAN}
Chúng em xin cam đoan rằng toàn bộ nội dung trong báo cáo Bài tập lớn này là kết quả nghiên cứu, tự cài đặt thuật toán và thực nghiệm độc lập của tập thể nhóm 5 sinh viên dưới sự hướng dẫn của TS. Phan Thanh Toàn.

Các thuật toán cốt lõi trong đề tài (lọc trung vị, phân ngưỡng Otsu tối ưu, các phép biến đổi hình thái học, trích xuất đường biên và đặc trưng hình học/cường độ) hoàn toàn được nhóm tự lập trình bằng thư viện ma trận \texttt{NumPy} từ số 0, tuyệt đối không sử dụng các hàm đen (black-box) của thư viện ngoài để thay thế thuật toán chính. Mọi số liệu thực nghiệm trên 5,630 ảnh SAR Sentinel-1 là trung thực và có thể tái lập từ mã nguồn trong kho lưu trữ Git của nhóm.

\begin{flushright}
    \textit{Hà Nội, ngày 17 tháng 10 năm 2026}\\
    \textbf{Đại diện nhóm sinh viên}\\
    \textit{(Ký và ghi rõ họ tên)}\\[1.5cm]
    \textbf{[Họ và tên Trưởng nhóm - TV1]}
\end{flushright}

\newpage
\section*{LỜI CẢM ƠN}
Nhóm xin gửi lời cảm ơn chân thành và sâu sắc nhất tới \textbf{TS. Phan Thanh Toàn} --- Giảng viên phụ trách học phần Xử lý ảnh, Khoa Công nghệ Thông tin, Trường Đại học Giao thông Vận tải. Thầy đã truyền đạt những kiến thức nền tảng vô cùng sâu sắc về cải thiện ảnh, xử lý hình thái học và phân đoạn ảnh, đồng thời luôn định hướng nhóm tiếp cận các bài toán thực tiễn gắn liền với chủ quyền biển đảo và bảo vệ môi trường biển Việt Nam.

Nhóm cũng xin cảm ơn tổ chức CSIRO Data61 và Cơ quan Vũ trụ Châu Âu (ESA) đã cung cấp công khai tập dữ liệu ảnh viễn thám radar Sentinel-1 SAR chất lượng cao làm nền tảng cho các nghiên cứu và thực nghiệm trong đề tài.

\newpage
% ==============================================================================
% MỤC LỤC & DANH MỤC TỰ ĐỘNG
% ==============================================================================
\tableofcontents
\newpage
\listoffigures
\newpage
\listoftables

\newpage
\section*{DANH MỤC TỪ VIẾT TẮT VÀ KÝ HIỆU TOÁN HỌC}
\addcontentsline{toc}{chapter}{DANH MỤC TỪ VIẾT TẮT VÀ KÝ HIỆU}

\begin{table}[H]
\centering
\begin{tabular}{ll}
\toprule
\textbf{Ký hiệu / Viết tắt} & \textbf{Ý nghĩa đầy đủ} \\
\midrule
\textbf{SAR} & Synthetic Aperture Radar (Radar khẩu độ tổng hợp) \\
\textbf{CDIO} & Conceive -- Design -- Implement -- Operate (Chuẩn đào tạo kỹ thuật) \\
\textbf{BTL} & Bài tập lớn môn học \\
\textbf{UTC} & University of Transport and Communications (ĐH Giao thông Vận tải) \\
\textbf{Otsu} & Thuật toán phân ngưỡng tối ưu dựa trên phương sai liên lớp \\
\textbf{CLAHE} & Contrast Limited Adaptive Histogram Equalization \\
\textbf{AIS} & Automatic Identification System (Hệ thống tự động nhận dạng tàu biển) \\
$I(x,y)$ & Cường độ mức xám tại tọa độ $(x, y)$ \\
$\mu_T$ & Mức xám trung bình toàn cục của ảnh \\
$\omega(k)$ & Xác suất tích lũy của các mức xám từ 0 đến $k$ \\
$\mu(k)$ & Giá trị mức xám trung bình tích lũy đến ngưỡng $k$ \\
$\sigma_B^2(k)$ & Phương sai giữa hai lớp (Inter-class variance) tại ngưỡng $k$ \\
$\sigma_T^2$ & Tổng phương sai toàn cục của ảnh (Total variance) \\
$\eta$ & Hệ số khả tách của thuật toán Otsu ($\eta = \sigma_B^2 / \sigma_T^2$) \\
$A \ominus B$ & Phép co (Erosion) tập nhị phân $A$ bởi phần tử cấu trúc $B$ \\
$A \oplus B$ & Phép dãn (Dilation) tập nhị phân $A$ bởi phần tử cấu trúc $B$ \\
$A \circ B$ & Phép mở hình thái học ($A \circ B = (A \ominus B) \oplus B$) \\
$A \bullet B$ & Phép đóng hình thái học ($A \bullet B = (A \oplus B) \ominus B$) \\
$\beta(A)$ & Trích biên hình thái học đơn pixel ($\beta(A) = A - (A \ominus B)$) \\
$M(x,y)$ & Độ lớn Gradient đạo hàm biên Sobel ($M = \sqrt{G_x^2 + G_y^2}$) \\
$\Delta \mu$ & Độ tương phản mức xám giữa vành đai biển lân cận và thân vết dầu \\
\bottomrule
\end{tabular}
\end{table}
